% experiments/exp2_chain_recipe/section8_recipe_revised.tex
%
% New Section 8 of draft.tex. Supersedes the four-piece recipe of
% Section 3 with a three-piece recipe distilled from the v14
% ablation series of Section 7 (section7_v14_findings.tex).
%
% This file is meant to be \input'd after section7_v14_findings.tex
% in the master draft. It does not modify draft.tex.

\section{The revised three-piece recipe: norm, no-warmup, and contrastive supervision}
\label{sec:revised_recipe}

The four-piece recipe of \S\ref{sec:recipe} and the three additional
training-dynamics pieces of \S\ref{sec:second_iteration} were each
motivated by a specific failure mode surfaced on the corresponding
corpus. With the v14 evidence of \S\ref{sec:third_iteration} in
hand, two of those pieces are now empirically retired and one is
replaced by a stricter architectural prerequisite. We close the
recipe at \emph{three} training-time pieces. They are: an
\textbf{input-RMSNorm on the extractor}, the \textbf{absence of any
supervised writer warmup}, and the combination of
\textbf{alpha-floor auxiliary loss} with \textbf{multi-negative
InfoNCE}. We restate each piece, list the recipe knobs that this
section retires, and give the canonical launch command.

\paragraph{Piece~1 (revised): input-RMSNorm on the extractor.}
The contextual extract source of \S\ref{sec:recipe}'s Piece~$1$
fed the per-token mid-stack hidden state $C_t =
\mathrm{hidden}_L(\text{input\_ids})$ directly into the extract
cross-attention. On Qwen3-0.6B at $L{=}14$ the residual stream's
RMS is $\approx 50$ in fp16, large enough that the
$W_K^{\text{ext}}$ backward gradients are
$\sim\!50\times$ what the writer's optimisation budget can absorb
in a single Adam step at $\eta_{\text{m}}{=}2\!\times\!10^{-4}$. We
add a single \texttt{Qwen3RMSNorm} layer to $C_t$ before the
extract cross-attention is computed, gated by the new flag
\texttt{--memres\_extract\_input\_norm}. The norm parameters are
trainable and start at the Qwen3 default
($\gamma_i \equiv 1$, $\varepsilon{=}10^{-6}$). Empirically, this
reduces the per-step extract.W\_K backward gradient by a factor
of $\sim\!50$ and is a necessary precondition for any
frozen-backbone training: without it, the cells in
Table~\ref{tab:v14_grid} that have warmup off (e.g.\ a no-warmup,
no-norm sibling we ran briefly off-table) produce
$W_K^{\text{ext}}$ updates large enough to send the writer
through a reset cascade within $\sim\!50$ steps. With it,
the writer optimises smoothly throughout joint training.

\paragraph{Piece~2 (negative): no supervised writer warmup.}
\S\ref{sec:third_iteration} establishes the empirical case against
\S\ref{sec:second_iteration}'s Piece~$7$. The revised recipe
explicitly turns the warmup off:
\texttt{--writer\_warmup\_steps 0
--writer\_warmup\_router\_bias 0.0
--writer\_warmup\_anneal\_steps 0}. The router starts at
$\texttt{mem\_bias\_init}{=}0$ and
$\texttt{recent\_bias\_init}{=}0$ --- a strictly weaker init than
the soft $\pm 4$ of \S\ref{sec:recipe}'s Piece~$4$. The motivation
for the $\pm 4$ soft init was to keep the bf16 backward gradient
through the routing softmax inside the format's representable range
under \emph{any} contrastive ramp; with the writer warmup retired,
the only term that has to backpropagate through the routing softmax
is the joint LM/InfoNCE/floor objective, and we observe in v14k
that an unbiased init produces a usable gradient on
$\texttt{mem\_bias}$ from step $1$. We retain
\S\ref{sec:negative_result}'s analytic argument against the strict
$\pm 32$ init unchanged --- that argument is about the $\pm 32$
saturation, not about whether $\pm 4$ is a numerical floor.

\paragraph{Piece~3: alpha-floor auxiliary loss $+$ multi-negative InfoNCE.}
The two remaining surviving pieces from
\S\ref{sec:second_iteration} are the $\amem$ load-balance auxiliary
($\mathcal{L}_{\text{floor}}$ with $\tau{=}0.05$,
$\lambda_{\text{floor}}{=}0.01$) and the multi-negative InfoNCE
contrastive on callback tokens
($\lambda_{\text{NCE}}{=}0.5$, $B{=}4$ batch-wise negatives,
temperature $\mathcal{T}{=}1.0$). They are unchanged from
\S\ref{sec:second_iteration}'s definitions. The combination provides
the discriminative signal the router needs to recruit memory beyond
plain NLL: the alpha-floor prevents MoE-style expert collapse of
the memory source under joint training, and InfoNCE makes the
writer's $M_{\text{new}}$ chain-discriminative on the callback
sessions where evidence is concentrated. Critically, on D4v2 these
two pieces are jointly insufficient \emph{without} the input-RMSNorm
of Piece~$1$ (the extract gradients run away within $\sim\!50$
steps) and jointly insufficient \emph{with} the writer warmup of
\S\ref{sec:second_iteration}'s Piece~$7$ (the post-anneal recent-bias
lock-in dominates). They are sufficient in v14k, where Piece~$1$ is
on, Piece~$7$ is off, and the alpha-floor saturates the inequality
near the floor as designed: $\amem^{\text{mean}}$ at step $3000$
sits at $0.052 \pm 0.003$ across the routed sublayers, just above
$\tau$, and per-sublayer max $\amem^{\text{max}}$ on
evidence-bearing tokens reaches $0.31$ on the most evidence-rich
sublayer, consistent with the writer concentrating its readout
when called for.

\paragraph{What is retired.}
Three pieces from prior iterations no longer appear in the recipe.
(i) The supervised writer warmup
(\texttt{--writer\_warmup\_steps},
\texttt{--writer\_warmup\_router\_bias},
\texttt{--writer\_warmup\_anneal\_steps}) is empirically
counterproductive on D4v2 and is set to zero. (ii) The soft $\pm 4$
parity router init of \S\ref{sec:recipe}'s Piece~$4$ is replaced by
the unbiased $\pm 0$ init; the analytic argument against
\emph{strict} $\pm 32$ in \S\ref{sec:negative_result} is unchanged,
but the soft $\pm 4$ floor was an over-correction once the writer
warmup is removed. (iii) The original Piece~$1$
(``contextual extract source'') is folded into the revised Piece~$1$
(``contextual extract source plus input-RMSNorm''): the contextual
hidden state is still the source, but the missing normalisation was
load-bearing. Pieces~$2$ and~$3$ of \S\ref{sec:recipe}
(mixed conversational corpus, negative-chain margin loss) are
\emph{not} retired but are folded into the corpus build and the
existing \texttt{--neg\_chain\_*} flags respectively, and are
unchanged.

\paragraph{Canonical launch command.}
The full revised recipe is the v14k launch command, with the v15a
replicate as the immediate target:

{\small\begin{verbatim}
python -u src/train_chain.py \
    --preset qwen3-0.6b-large \
    --memres_mode attention_parity \
    --memres_extract_source hidden_14 \
    --memres_extract_input_norm \
    --memres_writer_kind slot_attention \
    --memres_judge_qk_layernorm \
    --router_mem_bias_init 0   --router_recent_bias_init 0 \
    --writer_warmup_steps 0   --writer_warmup_router_bias 0.0 \
    --writer_warmup_anneal_steps 0 \
    --alpha_mem_floor_aux_weight 0.01 --alpha_mem_floor_target 0.05 \
    --contrastive_infonce_weight 0.5 --contrastive_infonce_callback_only \
    --train_chains paper_artifacts/chains/d4v2_train_s512.pt \
    --eval_chains  paper_artifacts/chains/d4v2_val_s512.pt \
    --window_k 8 --batch_size 4 --grad_accum 4 \
    --lr 2e-4 --lr_backbone 0.0 \
    --steps 3000 --warmup 200 --eval_every 200 --save_every 500 \
    --neg_chain_weight 0.5 --neg_chain_initial_weight 0.05 \
    --neg_chain_warmup_steps 1000 --neg_chain_margin 0.05 \
    --memory_dropout 0.10 --context_dropout 0.30 --carry_state \
    --burn_in_max 8 --burn_in_resample --mask_padding_loss \
    --score_tail_frac 1.0 --eval_window 4 --eval_n_chains 32 \
    --gradient_checkpointing --save_best_metric evidence_lift \
    --run_name chain_v15a_threepiece_d4v2
\end{verbatim}}

\noindent The diff from \S\ref{sec:recipe}'s
\texttt{chain\_v5\_softhidden14\_msc} command is six flags: the
addition of \texttt{--memres\_extract\_input\_norm} and the four
\texttt{--*\_qk\_layernorm} / \texttt{--memres\_writer\_kind}
architectural flags from \S\ref{sec:second_iteration}, and the
zeroing of all three \texttt{--writer\_warmup\_*} flags. The
\texttt{--lr\_backbone 0.0} is the frozen-backbone setting used in
the v14 series and in v15a; the v15b cell (Table~\ref{tab:v15_pending})
restores \texttt{--lr\_backbone 2e-5} to test joint training.

\begin{table}[t]
\centering
\caption{Summary of the revised three-piece recipe and the
corresponding \texttt{train\_chain.py} flags. Pieces~$2$
(contextual extract source) and~$3$ (mixed conversational corpus
$+$ negative-chain margin) of \S\ref{sec:recipe} are folded into
this recipe via the corpus build and the
\texttt{--neg\_chain\_*} flags respectively, and are unchanged.}
\label{tab:revised_recipe}
\small
\begin{tabular}{ll}
\toprule
piece & key flags \\
\midrule
1. extractor input-RMSNorm   & \texttt{--memres\_extract\_input\_norm} \\
\midrule
2. no writer warmup          & \texttt{--writer\_warmup\_steps 0} \\
                             & \texttt{--writer\_warmup\_router\_bias 0.0} \\
                             & \texttt{--writer\_warmup\_anneal\_steps 0} \\
                             & \texttt{--router\_mem\_bias\_init 0} \\
\midrule
3. alpha-floor $+$ InfoNCE   & \texttt{--alpha\_mem\_floor\_aux\_weight 0.01} \\
                             & \texttt{--alpha\_mem\_floor\_target 0.05} \\
                             & \texttt{--contrastive\_infonce\_weight 0.5} \\
                             & \texttt{--contrastive\_infonce\_callback\_only} \\
\bottomrule
\end{tabular}
\end{table}

\paragraph{Headline at scale (pending).}
The v15e and v15f cells port the revised recipe to Qwen3-1.7B; v15e
keeps the backbone frozen as in v14k, and v15f tests joint training
at $\eta_{\text{b}}{=}2\!\times\!10^{-5}$. Together with the v15a
$0.6$B replicate of v14k they form the proposed headline of the
revised manuscript: a $2\!\times\!2$ over
$\{\text{0.6B}, \text{1.7B}\} \times \{\text{frozen}, \text{joint}\}$
with the three-piece recipe held fixed. Table~\ref{tab:revised_headline}
gives the placeholder structure, with all entries gated on the
v15 runs currently in flight.

\begin{table}[t]
\centering
\caption{Planned headline of the revised three-piece recipe across
backbone scale and freeze regime. All cells use the
\S\ref{sec:revised_recipe} recipe (input-RMSNorm, no warmup,
alpha-floor $+$ InfoNCE) with the canonical launch command above.
The v15e and v15f rows (in bold) are the 1.7B headline numbers
gated on the revised manuscript. Metrics on the persisted best
checkpoint of each cell.}
\label{tab:revised_headline}
\small
\begin{tabular}{lcccc}
\toprule
cell & backbone $\times$ regime & D4v2 evidence\_lift & D4v2 $\dsh$ [CI] & LoCoMo $\dsh$ [CI] \\
\midrule
v15a & 0.6B, frozen      & \todo{} & \todo{} & \todo{} \\
v15b & 0.6B, joint       & \todo{} & \todo{} & \todo{} \\
\textbf{v15e} & \textbf{1.7B, frozen} & \todo{} & \todo{} & \todo{} \\
\textbf{v15f} & \textbf{1.7B, joint}  & \todo{} & \todo{} & \todo{} \\
\bottomrule
\end{tabular}
\end{table}

\paragraph{Reading.}
The revised recipe replaces the seven pieces of \S\S\ref{sec:recipe},
\ref{sec:second_iteration} with three. Two of the original pieces
(corpus mix, negative-chain margin) survive unchanged but are
folded into infrastructure rather than highlighted as recipe
contributions, since neither was a knob in the v14 ablation. Two
were retired by the v14 evidence: the supervised writer warmup,
which is counterproductive on D4v2, and the soft $\pm 4$ router
init, which was an over-correction necessitated only by the
warmup phase. One was strengthened: the contextual extract source
of \S\ref{sec:recipe}'s Piece~$1$ now requires an input-RMSNorm to
be operational under a frozen backbone. The headline empirical
claim of the revised manuscript is that this three-piece recipe is
the minimal recipe that produces a sustained positive
evidence\_lift on a controlled callback corpus at $0.6$B; the v15
series tests how far that claim ports along the freeze and scale
axes.
